\documentclass[10pt]{article}
\usepackage[T1]{fontenc}
\usepackage[utf8]{inputenc}
\usepackage[french]{babel}
\usepackage{amsmath,amssymb,mathrsfs,mathpazo}
\usepackage{geometry}
\geometry{a4paper,margin=2cm}
\usepackage{graphicx,array,tabularx,tabularray,colortbl,multirow,multicol,booktabs,enumitem,pifont,ulem,fancybox,diagbox,cancel}
\usepackage{pstricks,pstricks-add,pst-plot,pst-tree,pst-node,pst-text,pst-func,pst-math,pst-bspline,pst-3dplot,pst-eucl,pst-solides3d}
\let\pscalculate\undefined
\usepackage{pst-bezier}
\usepackage{tikz,pgfplots,tkz-tab}
\usetikzlibrary{arrows,arrows.meta,patterns,calc,positioning,trees,decorations.markings,shapes.geometric,shapes.misc,shapes.arrows,intersections,decorations.pathmorphing,decorations.pathreplacing,backgrounds,shadows,babel}
\pgfplotsset{compat=1.18}
\usepgfplotslibrary{fillbetween}
\usepackage[np]{numprint}
\usepackage{listings,makecell,pas-tableur,twoopt}
\let\barre\undefined
\usepackage[pstricks]{tablvar}
\definecolor{darkgreen}{RGB}{0,100,0}
\providecommand{\red}{\color{red}}
\providecommand{\blue}{\color{blue}}
\providecommand{\green}{\color{green}}
\providecommand{\black}{\color{black}}
\providecommand{\white}{\color{white}}
\providecommand{\gray}{\color{gray}}
\providecommand{\cyan}{\color{cyan}}
\providecommand{\magenta}{\color{magenta}}
\providecommand{\yellow}{\color{yellow}}
\providecommand{\violet}{\color{violet}}
\providecommand{\orange}{\color{orange}}
\providecommand{\euro}{\texteuro}
\providecommand{\EUR}{\texteuro}
\providecommand{\ieme}{\textsuperscript{e}}
\providecommand{\degres}{\ensuremath{^\circ}}
\providecommand{\degree}{\ensuremath{^\circ}}
\providecommand{\vv}[1]{\overrightarrow{#1}}
\providecommand{\Coord}[1]{#1}
\providecommand{\decosix}{}
\providecommand{\starredbullet}{\ensuremath{\bullet}}
\pagestyle{empty}
\setlength{\parindent}{0pt}
\usepackage{adjustbox}\DeclareUnicodeCharacter{2236}{:}\DeclareUnicodeCharacter{2190}{\ensuremath{\leftarrow}}
\providecommand{\R}{\mathbb{R}}
\providecommand{\N}{\mathbb{N}}
\providecommand{\Z}{\mathbb{Z}}
\providecommand{\Q}{\mathbb{Q}}
\providecommand{\D}{\mathbb{D}}
\providecommand{\C}{\mathbb{C}}
\providecommand{\e}{{\mathrm{e}}}
\providecommand{\vect}[1]{\overrightarrow{#1}}
\providecommand{\vectt}[1]{\overrightarrow{\mathrm{#1}}}
\providecommand{\vvt}[1]{\overrightarrow{\mathrm{#1}}}
\providecommand{\np}[2][]{{#2~#1}}
\providecommand{\vv}[1]{{\overrightarrow{#1}}}
\renewcommand{\d}{{\mathrm{d}}}
\providecommand{\degres}{{^\circ}}
\providecommand{\degre}{{^\circ}}
\providecommand{\blue}{\color{blue}}
\providecommand{\red}{\color{red}}
\providecommand{\enskip}{\quad}
\providecommand{\llbracket}{\lbrack\!\lbrack}
\providecommand{\rrbracket}{\rbrack\!\rbrack}
\providecommand{\euro}{{\text{€}}}
\providecommand{\no}{{\text{n°}}}
\providecommand{\interval}[2]{{[#1;#2]}}
\providecommand{\intervFO}[2]{{[#1;#2[}}
\providecommand{\dotfill}{\ldots}

\providecommand{\ds}{\displaystyle}
\providecommand{\textding}[1]{#1}
\providecommand{\Oij}{\(\mathrm{O};\vec{\imath},\vec{\jmath}\)}
\providecommand{\Oijk}{\(\mathrm{O};\vec{\imath},\vec{\jmath},\vec{k}\)}
\providecommand{\Ouv}{\(\mathrm{O};\vec{u},\vec{v}\)}
\providecommand{\intervFF}[2]{\ensuremath{[#1;#2]}}
\renewcommand{\euro}{\ifmmode\text{\texteuro}\else\texteuro\fi}
\DeclareUnicodeCharacter{00B0}{\ifmmode^{\circ}\else\textdegree\fi}
\DeclareUnicodeCharacter{20AC}{\ifmmode\text{\texteuro}\else\texteuro\fi}
% Énoncés : Éducation nationale. Transcription et corrigés : APMEP, Denis Vergès et contributeurs. Découpage et classement : CoopMaths, Sébastien Lozano, Éric Elter et contributeurs. Adaptation de la présentation : Maths Explorer.
% CC BY-SA : https://coopmaths.fr/www/about/
% https://coopmaths.fr/alea/static/e3c/2020/tex/e3c_2020_00_sujet24_2_cor.tex
\begin{document}
\noindent\textit{APMEP / CoopMaths. CC BY-SA. Présentation adaptée par Maths Explorer.}\par\medskip
\subsection*{Question 1}

La tangente en $A$ est horizontale, donc $f'(-1) = 0$.

\subsection*{Question 2}

On lit sur le graphe que $f'(x) \leq 0$ sur $[-4 ; -1]$ et sur $[1,5 ; 2]$.

\subsection*{Question 3}

On admet que la fonction $f$ est définie sur $[-4 ; 2]$ \\
par $f(x) = (-x^2 + 2,5x - 1)e^x$.\\
La fonction $f$ est un produit de fonctions dérivables sur $[-4 ; 2]$, donc sur cet intervalle :
\begin{align*}
	f'(x) &= (-2x + 2,5)e^x + (-x^2 + 2,5x - 1)e^x \\
	&= e^x \left( -2x + 2,5 - x^2 + 2,5x - 1 \right) \\
	&= e^x (-x^2 + 0,5x + 1,5).
\end{align*}

\subsection*{Question 4}

Le signe de $f'(x)$ dépend du signe du trinôme $-x^2 + 0,5x + 1,5$.\\
Le discriminant de ce trinôme est :
\begin{align*}
	\Delta &= 0,5^2 + 4 \times 1,5 \\
	&= 0,25 + 6 = 6,25 > 0.
\end{align*}
Ce trinôme a donc deux racines :
\begin{align*}
	x_1 &= \frac{-0,5 + 2,5}{-2} = -1, \\
	x_2 &= \frac{-0,5 - 2,5}{-2} = 1,5.
\end{align*}
Le trinôme est négatif sur $[-4 ; -1[$ et sur $[1,5 ; 2]$ et positif sur $[-1 ; 1,5]$.

Ainsi :
\begin{itemize}
	\item $f'(x) < 0$ sur $[-4 ; -1[$ et sur $[1,5 ; 2]$, donc la fonction $f$ est décroissante sur ces intervalles ;
	\item $f'(x) > 0$ sur $[-1 ; 1,5]$, donc la fonction $f$ est croissante sur cet intervalle ;
	\item $f'(-1) = 0$ et $f'(1,5) = 0$, donc $f(-1)$ et $f(1,5)$ sont les extremums de $f$ sur $[-4 ; 2]$.
\end{itemize}

\subsection*{Question 5}

D'après la question précédente, on a le tableau de variations suivant :

\begin{center}
\begin{tikzpicture}
	\tkzTabInit[lgt=3, espcl=3]
	{$x$/1, $f'(x)$/1, $f(x)$/2}
	{$-4$, $-1$, $1.5$, $2$}
	\tkzTabLine{,-,z,+,z,-}
	\tkzTabVar{+/$-27\e^{-4}$, -/$-4.5\e^{-1}$, +/$0.5\e^{4}$, -/$0$}
\end{tikzpicture}
\end{center}

\end{document}
